\documentclass[11pt,a4paper]{article}
\usepackage{amssymb}

\usepackage[margin=2.5cm]{geometry}

\usepackage{helvet}
\renewcommand{\familydefault}{\sfdefault}
\usepackage[T1]{fontenc}
\usepackage{microtype}
\usepackage{enumitem}
\usepackage{hyperref}
\usepackage{titlesec}
\usepackage{parskip}

\titleformat{\section}{\large\bfseries}{\thesection}{1em}{}
\titleformat{\subsection}{\normalsize\bfseries}{\thesubsection}{1em}{}

\hypersetup{
    colorlinks=true,
    linkcolor=black,
    urlcolor=black
}

\title{\textbf{Grim Frontiers Web Application}\\
\large Beta Scope Proposal}
\author{}
\date{\today}

\begin{document}

\maketitle

\section*{Objective}

Deliver a functional web-based application that:

\begin{itemize}[leftmargin=2em]
    \item Allows multiplayer sessions (Marshal + Players)
    \item Handles the core card mechanics (Ashcan core rules, pp. 13--15)
    \item Assists scene resolution
    \item Supports solo play for rules learning
    \item Is stable, usable, and technically solid
\end{itemize}

This beta is a \textbf{rules engine and table assistant}, not a full narrative platform.

\section{Multiplayer (Beta Level)}

\subsection*{Included}

\begin{itemize}[leftmargin=2em]
    \item Create a game room (unique code)
    \item Join via room code (no account system required)
    \item Role assignment:
        \begin{itemize}
            \item 1 Marshal
            \item N Players
        \end{itemize}
    \item Each participant sees:
        \begin{itemize}
            \item Their own deck
            \item Their own hand
        \end{itemize}
    \item Shared table view:
        \begin{itemize}
            \item Scene state
            \item Draw pile size
            \item Discard pile
            \item Event log
        \end{itemize}
\end{itemize}

\subsection*{Not Included in Beta}

\begin{itemize}[leftmargin=2em]
    \item Persistent user accounts
    \item Campaign storage across sessions
    \item Voice or video integration
    \item In-app chat (external communication assumed)
    \item Native mobile application
\end{itemize}

\section{Core Mechanics Implemented}

\subsection{Deck Handling}

\begin{itemize}[leftmargin=2em]
    \item Standard 52-card deck (+ Jokers as per rules)
    \item Shuffle
    \item Draw
    \item Discard
    \item Reshuffle when required
    \item Visible discard pile
    \item Automatic Joker handling (if required by rules)
\end{itemize}

\subsection{Resolution System}

\begin{itemize}[leftmargin=2em]
    \item Action declaration
    \item Card draw resolution
    \item Total calculation (Blackjack-style logic where applicable)
    \item Success / Partial / Failure determination
    \item Damage assignment
    \item Condition tracking
\end{itemize}

\subsection{Dirty Business (Non-Combat Resolution)}

\begin{itemize}[leftmargin=2em]
    \item Structured resolution flow
    \item Marshal-triggered draws where required
    \item Automatic outcome logging
\end{itemize}

\section{Scene Assistant}

The application supports structured scene flow:

\begin{itemize}[leftmargin=2em]
    \item Start Scene
    \item Track active participants
    \item Track damage and conditions
    \item Resolve action
    \item Log outcomes
\end{itemize}

All actions generate a timestamped event log.

\textbf{Example Log Output:}

{\ttfamily
[Turn 3] Player A draws 9$\spadesuit$ \newline
[Turn 3] Total = 18 → Success \newline
[Turn 3] Marshal draws Q$\diamondsuit$
}

Session logs can be exported (JSON or text).

\section{Solo Mode}

Purpose: Learn mechanics and test rules.

\begin{itemize}[leftmargin=2em]
    \item Single player vs automated Marshal deck logic
    \item App performs Marshal draws automatically
    \item No narrative generation
    \item Simple scene loop
    \item Reset functionality
\end{itemize}

This is rules automation, not AI storytelling.

\section{User Interface Philosophy}

\begin{itemize}[leftmargin=2em]
    \item Clean
    \item Functional
    \item Minimal
    \item No animations required
    \item Cards displayed as:
        \begin{itemize}
            \item Simple symbols (e.g., Q$\spadesuit$), or
            \item Static card images (optional)
        \end{itemize}
\end{itemize}

Primary focus: clarity and stability.

\section{Technical Approach (Beta)}

\begin{itemize}[leftmargin=2em]
    \item Backend: Python (FastAPI)
    \item Real-time synchronization: WebSockets
    \item Storage: SQLite (beta-level persistence)
    \item Frontend: Minimal React interface
    \item Hosting: VPS or staging server
\end{itemize}

\section{Explicit Beta Limitations}

This beta:

\begin{itemize}[leftmargin=2em]
    \item Is not final product quality
    \item Is not optimized for mobile
    \item Is not load-tested
    \item Is not a publishing platform
    \item Is not feature-complete
\end{itemize}

It is a proof-of-concept to validate:

\begin{itemize}[leftmargin=2em]
    \item Gameplay flow
    \item Multiplayer feasibility
    \item Solo learning mode
    \item Technical viability
\end{itemize}

\section{Proposed Timeline}

\textbf{Week 1:} Game engine module (backend logic only)

\textbf{Week 2:} Room system + WebSocket synchronization

\textbf{Week 3:} Core resolution + damage + logging

\textbf{Week 4:} Solo mode automation

\textbf{Week 5:} Stabilization and minimal polish

\bigskip

\textbf{Target:} End of May – Closed Beta Demonstration

\section*{Scope Discipline}

Any of the following would move this from ``beta'' to ``startup'' and are intentionally excluded:

\begin{itemize}[leftmargin=2em]
    \item Campaign persistence
    \item Character sheet system
    \item XP progression
    \item Map system
    \item AI narration
    \item Card animations
    \item Marketplace or store integration
\end{itemize}

\bigskip

\noindent
\textit{The purpose of this beta is to validate mechanics, not to finalize the platform.}

\end{document}
